\scanpage{10}
l'hom. naturel
\[
\underline{\operatorname{Hom}}(P,\Omega)_s\otimes\mathcal O_{S,s}/(f_1\ldots f_n)\longrightarrow
\underline{\operatorname{Hom}}_{\mathcal O_{S,s}}\bigl(P_s,\Omega_s/(f_1\ldots f_{\unsure{s|n}})\Omega_{\unsure{s}}\bigr)
\]
soit un isom.

\uline{Dém} On peut supposer $S$ affine \inserted{anneau $A$}, alors on a \deleted{résolution} présentation
\[
L_1\mathbin{-}L_0\longrightarrow P
\]
avec $L_1,L_0$ \unsure{libre|libres} de t.f., donc $\underline{\operatorname{Hom}}(P,Q)$ est le $H^0$ du complexe
\[
K^{\cdot}:\qquad 0-\overset{\text{degré }0}{\check L_0\otimes P}-\overset{\text{degré }1}{\check L_1\otimes P}-0
\]
tandis que $\operatorname{Hom}(P,Q\otimes M)$ est le $H^0$ du complexe $K^{\cdot}\otimes M$. Un critère pour que l'hom naturel
\[
H^0(K^{\cdot})\otimes M^{\unsure{\cdot}}-\ H^0(K^{\cdot}\otimes M)
\]
soit bijectif est qu'on ait
\[
\operatorname{Tor}_1(M,H^1(K^{\cdot}))=\operatorname{Tor}_1(M,Z^1(K^{\cdot}))=0
\]
Or $H^1(K)=R$ et $Z^1(K)=S$ sont des modules de t.f. sur $A$, auxquels on peut appliquer le \deleted{corollaire} lemme 6, b), en faisant $M=\mathcal O_{S,s}/\unsure{(f_1\ldots f_n)|f_1\ldots f_n}$.

On gagne.

% archive sheet number: 9
